\documentclass[10pt,a4paper]{amsart}
\usepackage[margin=19mm]{geometry}
\usepackage{amsmath,amssymb,mathtools}
\usepackage[scr=rsfs,cal=euler]{mathalfa}
\usepackage{xfrac}
\usepackage[unicode,hidelinks,bookmarks=false]{hyperref}
\newcommand{\ie}{i.e.\ }
\newcommand{\cf}{cf.\ }
\newcommand{\resp}{respectively\ }
\newcommand{\Subsection}{\subsection}
\providecommand{\nobreakdash}{\nobreak\hbox{-}}
\setlength{\emergencystretch}{2em}
\multlinegap0pt
\usepackage{bm}
\usepackage{bbm}
\usepackage{dsfont}
\usepackage{xspace}
\usepackage{cancel}
\usepackage{mathtools}
\usepackage{cleveref}
\usepackage{amsthm}
\makeatletter
\@ifundefined{theorem}{\newtheorem{theorem}{Theorem}}{}
\@ifundefined{lemma}{\newtheorem{lemma}[theorem]{Lemma}}{}
\@ifundefined{proposition}{\newtheorem{proposition}[theorem]{Proposition}}{}
\makeatother
\newcommand{\ELR}{\mathrm{ELR}}
\newcommand{\R}{\mathbb{R}}
\title[Bounded Treewidth and Complete Monotonicity for Scott-Sokal ]{Bounded Treewidth and Complete Monotonicity for Scott-Sokal Spanning-Tree Polynomials}
\author{Domingos S. P. Salazar}
\date{Source: arXiv:2606.26275v1 (CC BY 4.0)}
\begin{document}
\maketitle

\noindent\emph{Source and licence.} Bounded Treewidth and Complete Monotonicity for Scott-Sokal Spanning-Tree Polynomials. Version arXiv:2606.26275v1 (CC BY 4.0), distributed under the Creative Commons Attribution 4.0 International licence, as recorded in the arXiv licence metadata for this exact version and shown on the abstract page https://arxiv.org/abs/2606.26275v1. Full rights notice: licenses/arxiv-2606.26275-CC-BY-4.0.txt.\par\smallskip

\noindent\textit{Editorial scope.} Counted unit: the Proposition whose source label is \texttt{prop:simplicial-extension-elr} in the article (source0.tex lines 533--540, source label \texttt{prop:simplicial-extension-elr}), delivered together with the article's own complete original proof of it (source0.tex lines 542--612, one \texttt{proof} environment of 71 lines). No mathematics is written, repaired or re-derived: statement and proof are the article's own text line for line. Required context retained uncounted: lem:starmesh (lines 350--363); lem:gaussian-star (lines 391--429). Every definition, display and label of those lines is retained unchanged. Omitted from this package and not redistributed: 1--349, 364--390, 430--532, 541--541, 613--1112. Nothing inside a retained excerpt is omitted or altered: every retained body is delivered exactly as the article's lines are written, so no omission receipt of any kind accompanies this package. Citations: the retained excerpts carry no citation command at all, so no bibliography entry of the article is transcribed and no reference is redistributed. Reference adaptation: none was needed and none was made; every reference inside the delivered unit resolves inside it, so no unresolved reference or citation is produced. Printed slips kept literal: no printed word, symbol, hypothesis or number of the counted statement or of its proof was altered. Layout and class: the article's own class is \texttt{elsarticle} with packages including fontenc, inputenc, amsmath, amssymb, amsthm, mathtools, microtype, tikz, hyperref, cleveref; the wrapper is the lane's pinned \texttt{amsart} editorial wrapper at 10pt on A4, which restates the theorem-like environments the retained text needs and adds the article's own macro definitions that the retained text needs (\texttt{\textbackslash ELR}, \texttt{\textbackslash R}), with extra wrapper packages (bm, bbm, dsfont, xspace, cancel, mathtools, cleveref). The delivered numbering is the unit's own. Assets: no retained span contains a figure environment, a graphics-inclusion command, a PDF-page inclusion, a table or a TikZ picture; no image file is copied into this package, the package declares no asset dependency and no image file is redistributed with it. Licence: version arXiv:2606.26275v1 of this article is distributed under the Creative Commons Attribution 4.0 International licence, as recorded in the arXiv licence metadata for this exact version and shown on the abstract page https://arxiv.org/abs/2606.26275v1. A full rights notice is delivered at licenses/arxiv-2606.26275-CC-BY-4.0.txt. Characters: every retained excerpt is pure ASCII, so the delivered engine, XeLaTeX with its default Latin Modern text font, reports no missing character.
\begin{lemma}[Star--mesh elimination]
\label{lem:starmesh}
Let \(H=(V,E_H)\) be a connected network and let \(S\subseteq V\) have size \(d\ge1\).  Form \(G\) by adjoining a new vertex \(v\) adjacent exactly to the vertices of \(S\), with conductances \(u_i>0\) for \(i\in S\).  Put
\begin{equation}
        U=\sum_{i\in S}u_i,
        \qquad
        p_{ij}(u)=\frac{u_i u_j}{U}
        \quad (i<j,\ i,j\in S).
\end{equation}
Let \(x+p(u)\) denote the edge-weight vector of the network obtained from \(H\) by adding \(p_{ij}(u)\) to the conductance between \(i\) and \(j\) for every distinct \(i,j\in S\), creating that terminal edge if it was absent, and leaving all other edge weights unchanged.  Then
\begin{equation}
        T_G(x,u)=U\,T_H(x+p(u)).
\end{equation}
\end{lemma}

\begin{lemma}[Gaussian star kernel]
\label{lem:gaussian-star}
Let \(d\ge1\), let \(\xi_1,\ldots,\xi_d\) be points in a real Euclidean space, and put
\begin{equation}
        r_{ij}=\|\xi_i-\xi_j\|^2,
        \qquad
        U=\sum_{i=1}^d u_i,
        \qquad
        u_i>0.
\end{equation}
If
\begin{equation}
        \beta>\frac{d-1}{2},
\end{equation}
then
\begin{equation}
        K(u)
        =
        U^{-\beta}
        \exp\left(
        -\frac1U\sum_{1\le i<j\le d}u_i u_j r_{ij}
        \right)
\end{equation}
is a positive Laplace transform in \((u_1,\ldots,u_d)\).  More precisely, after embedding the original Euclidean space into a larger Euclidean space if necessary, let \(B\) be any \((d-1)\)-dimensional affine subspace containing \(\xi_1,\ldots,\xi_d\).  Then
\begin{equation}
\begin{aligned}
        K(u)
        &=
        \frac{\pi^{-(d-1)/2}}{\Gamma(\beta-(d-1)/2)}
        \int_0^\infty\int_B
        t^{\beta-(d-1)/2-1}  \\
        &\quad\times
        \exp\left(
        -\sum_{i=1}^d u_i\bigl(t+\|z-\xi_i\|^2\bigr)
        \right)\,dz\,dt .
\end{aligned}
\end{equation}
For \(d=1\), \(B=\{\xi_1\}\) and the empty pair sum is zero.  In particular, for \(d=3\) the kernel used in the parallel star-lift proof is completely monotone for every \(\beta>1\).
\end{lemma}

\begin{proposition}[Simplicial extension preserves \(\ELR_\beta\)]
\label{prop:simplicial-extension-elr}
Let \(H=(V,E_H)\) be a connected graph having \(\ELR_\beta\).  Let \(S\subseteq V\) be a clique of size \(d\ge1\), and let \(G\) be obtained from \(H\) by adjoining a new vertex \(v\) adjacent exactly to the vertices of \(S\).  If
\begin{equation}
        \beta>\frac{d-1}{2},
\end{equation}
then \(G\) has \(\ELR_\beta\).
\end{proposition}

\begin{proof}
Let \(u_i=x_{vi}>0\) for \(i\in S\), and put \(U=\sum_{i\in S}u_i\).  By \cref{lem:starmesh},
\begin{equation}
        T_G(x,u)=U\,T_H(x+p(u)).
\end{equation}
Since \(H\) has \(\ELR_\beta\), there are \(N\) and a positive measure \(\mu\) on \((\R^N)^V\) such that
\begin{equation}
        T_H(y)^{-\beta}
        =
        \int
        \exp\left(
        -\sum_{ij\in E_H}y_{ij}\|\xi_i-\xi_j\|^2
        \right)\,d\mu(\xi).
\end{equation}
Substituting \(y=x+p(u)\) gives
\begin{equation}
\begin{aligned}
        T_G(x,u)^{-\beta}
        &=
        \int
        \exp\left(
        -\sum_{ij\in E_H}x_{ij}\|\xi_i-\xi_j\|^2
        \right)  \\
        &\quad\times
        U^{-\beta}
        \exp\left(
        -\frac1U
        \sum_{\substack{i<j\\ i,j\in S}}
        u_i u_j\|\xi_i-\xi_j\|^2
        \right)\,d\mu(\xi).
\end{aligned}
\end{equation}
For each fixed \(\xi\), the second line is the Gaussian star kernel of \cref{lem:gaussian-star}.

It remains only to record the measure-theoretic construction of the new configuration measure.  Put \(N_0=\max\{N,d-1\}\) and embed \(\R^N\) isometrically into \(\R^{N_0}\).  For each configuration \(\xi\), let \(a(\xi)=d^{-1}\sum_{i\in S}\xi_i\).  Apply Gram--Schmidt, with a fixed deterministic tie-breaking rule using the standard basis of \(\R^{N_0}\), to obtain measurably an orthonormal \((d-1)\)-frame \(e_1(\xi),\ldots,e_{d-1}(\xi)\) whose span contains the vectors \(\xi_i-a(\xi)\), \(i\in S\).  The deterministic tie-breaking is included only to make this a Borel choice of frame, so that the push-forward measure below is unambiguous.  Thus
\begin{equation}
        z_{\xi}(s)
        =
        a(\xi)+\sum_{\ell=1}^{d-1}s_\ell e_\ell(\xi),
        \qquad s\in\R^{d-1},
\end{equation}
parametrizes a \((d-1)\)-dimensional affine subspace containing the clique points.

Define a map to configurations in \((\R^{N_0+1})^{V\cup\{v\}}\) by
\begin{equation}
        \xi'_w=(\xi_w,0)\quad (w\in V),
        \qquad
        \xi'_v=(z_{\xi}(s),\sqrt t).
\end{equation}
Then, for \(i\in S\),
\begin{equation}
        \|\xi'_v-\xi'_i\|^2=t+\|z_{\xi}(s)-\xi_i\|^2,
\end{equation}
and old edge distances are unchanged.  Let \(\nu\) be the push-forward of
\begin{equation}
        d\mu(\xi)\,
        \frac{\pi^{-(d-1)/2}}{\Gamma(\beta-(d-1)/2)}
        t^{\beta-(d-1)/2-1}\,dt\,ds .
\end{equation}
All integrands are nonnegative, so Tonelli's theorem applies.  Substituting the Gaussian star-kernel formula gives
\begin{equation}
        T_G(x,u)^{-\beta}
        =
        \int
        \exp\left(
        -\sum_{ij\in E_H}x_{ij}\|\xi'_i-\xi'_j\|^2
        -\sum_{i\in S}u_i\|\xi'_v-\xi'_i\|^2
        \right)\,d\nu(\xi').
\end{equation}
This is precisely an \(\ELR_\beta\) representation for \(G\).
\end{proof}

\end{document}
